% Part: proof-theory
% Chapter: cut-elimination
% Section: introduction

\documentclass[../../../include/open-logic-section]{subfiles}

\begin{document}

\olfileid[tr]{pt}{cut}{int}

\olsection{Giriş}

\Cut{} kuralı şu kuraldır:
\[
\Axiom$\Gamma \fCenter \Delta, !A$
\Axiom$!A, \Pi \fCenter \Lambda$
\RightLabel{\Cut}
\BinaryInf$\Gamma, \Pi \fCenter \Delta, \Lambda$
\DisplayProof
\]
(!!{formula}~$!A$'ya \emph{kesme formülü} denir.)

\Cut{}, Gentzen'in özgün sekant hesabı~$\Log{LK}$'ye katılmıştı. Gentzen'in
\emph{Hauptsatz}ı (``ana teorem''), $\Log{LK}$ içinde !!{provable} olan her
sekantın \Cut{} kuralı kullanılmadan da !!{provable} olduğunu, yani
\Cut{}'ın $\Log{LK}$ ispatlarından ``giderilebildiğini'' söyler. $\Log{G1c}$
ile $\Log{G3c}$ sistemlerimiz \Cut{} içermez; fakat aynı soru onlar için de
sorulabilir: \Cut{}, \Cut{} ile birlikte $\Log{G1c}$'nin (ya da
$\Log{G3c}$'nin) öteki kurallarını kullanan !!{proof}s{}dan giderilebilir mi?
Yerleşik terminolojimizde bu, denk biçimde şöyle sorulabilir: \Cut{},
$\Log{G1c}$ (ya da $\Log{G3c}$) için kabul edilebilir bir kural mıdır?

Kesme-giderimi teoremi belki de ispat kuramının en önemli ve en ünlü
sonucudur. Öncelikle $\Log{G1c}$ ile $\Log{G3c}$'nin, ilk bakışta açıkça
gerekli olduğunu düşünebileceğiniz bir kurala ihtiyaç duymadığını gösterir.
Gerçek ispatları biçimselleştirmeye çalıştığınızda, önermelerin kullanımını
ele almanın bir yoluna ihtiyaç duyduğunuzu fark edersiniz. Matematikçiler
sonuçlar ispatlayıp başka sonuçların ispatlarında bunlara başvurmayı sever.
Dolayısıyla önce $L$ sonucunun (önermenin) bir ispatını $\Sequent L$
sekantının !!{proof}{}ı olarak biçimselleştirebilirsiniz. Ardından ispatında
$L$'yi kullanan ikinci $R$ sonucunun ispatını, $L \Sequent R$ sekantının
!!{proof}{}ı olarak biçimselleştirirsiniz. Yalnızca $\Sequent R$ sekantının
!!a{proof}{}ının, yani $R$'nin bir ispatının bulunduğunu nasıl bilirsiniz? İki
!!{proof}s{}ı birleştirmenin bir yoluna ihtiyaç vardır; \Cut{} kuralı bunu
yapmanızı sağlar.

$\Log{G1c}$ ile $\Log{G3c}$'nin tam olduğunu ve dolayısıyla ispatlanmasını
bekleyebileceğiniz bütün sekantları---bütün geçerli sekantları---ispatladığını
belirttik. Bu doğrudan ispatlanabilir. Fakat Gentzen yazarken yalnızca
aksiyomatik türetim sistemlerinin tam olduğu biliniyordu. Gentzen,
$\Log{LK}$'nin tam olduğunu göstermek için aksiyomatik sistemdeki
!!{derivation}s{}i $\Log{LK}$-!!{proof}s{}ına çeviren bir aktarım verdi. Bu
aktarım \Cut{}'ı esaslı biçimde kullanıyordu. Kesme-giderimi yalnızca klasik
mantık için önemli bir sonuç değildir: Başka pek çok mantığın sekant hesabı
vardır. Bazı mantıklarda sekant hesabının doğrudan tamlık ispatı zor ya da
olanaksızdır; dolayısıyla başka sistemlerden yapılan (\Cut{} kullanan)
aktarımlar tamlığı göstermenin tek yoludur. Bu durumda \Cut{}'ın gereksiz ve
\Cut{} içermeyen sistemin tam olduğunu göstermek için kesme-giderimini ispat
kuramsal yoldan ispatlamak gerekir.

Kesme-giderimi, bağımsız olarak ilginç başka sonuçlar elde etmek için de
uygulanabildiğinden önemlidir. Ele alacağımız iki örnek Craig aradeğerleme
önermesi ile Herbrand teoremidir.

Kesme-giderimini ispatlamak için kullanılan fikirler genellikle \Cut{}
kullanan !!a{proof}{}ın, onu kullanmayan bir ispata nasıl dönüştürüleceğini
göstermeyi içerir. Bu dönüştürmeler genellikle ``yereldir'': Bir !!a{proof}{}ın
bir parçasını (çoğu kez \Cut{} çıkarımıyla biten bir alt-!!{proof}{}ı) alıp
yerine aynı sekantın, \Cut{} çıkarımının çıkarıldığı ya da başka çıkarımlarla
değiştirildiği bir alt-!!{proof}{}ını koyarız. Böyle gerekçeler \Cut{} içeren
!!{proof}s{}ı, içermeyenlere dönüştüren adım adım algoritmalar verir. Bu,
kesme-giderimi yordamları uygulamalarda kullanıldığında daha fazla bilgi
sağlar. Örneğin aradeğerleme önermesinin ``$!A \lif !B$ geçerliyse bir
aradeğerleme vardır'' biçiminde model kuramsal ispatları bulunur (şimdilik
aradeğerlemenin ne olduğunu boş verin). Kesme-giderimini kullanan ispat
kuramsal bir gerekçe, $!A \lif !B$'nin !!a{proof}{}ını girdi alıp bir
aradeğerleme döndüren algoritma verir. Model kuramsal gerekçenin tersine,
ispat kuramsal gerekçe \emph{kurucudur}.

Kesme-giderimi teoremini ispatlamanın iki yaygın yolu vardır. Biri
(Gentzen'e uzanır), bir ispattaki en üst kesmelerin çıkarılabildiğini gösterir
ve hiç kesme kalmayana dek en üst kesmeleri çıkarır. Öteki (Schütte ve Tait'e
uzanır), !!a{proof} içindeki en karmaşık kesmelere odaklanır ve başka hiçbir
kesmenin daha yüksek derinlikte kesme !!{formula}{}ü bulunmaması anlamında
maksimal olan kesmeleri art arda çıkarır. Her iki yöntemin de üstün ve zayıf
yanları vardır. Bazı sistemlerde bir yaklaşım işlerken ötekini uygulamak zor,
hatta olanaksızdır. Bu nedenle iki yaklaşımı da açıklıyoruz. Kesme-giderimini
$\Log{G3c}$ için ispatlayacağız. Gerçekte \Cut{}'ın değil, \emph{bağlam
paylaşan kesme}~\CutCS{} kuralının giderilebildiğini göstereceğiz:
\[
\Axiom$\Gamma \fCenter \Delta, !A$
\Axiom$!A, \Gamma \fCenter \Delta$
\RightLabel{\Cut}
\BinaryInf$\Gamma \fCenter \Delta$
\DisplayProof
\]
Bir sekant hesabı zayıflatma ve büzülmeye (ya $\Log{G1c}$'de olduğu gibi
gerçek kurallar ya da $\Log{G3c}$'de olduğu gibi kabul edilebilir kurallar
olarak) izin veriyorsa \Cut{}, \CutCS{}'yi; \CutCS{} de \Cut{}'ı
taklit edebilir. İlk kesme-giderimi stratejisi \CutCS{} için daha kolay
açıklandığından ve ikinci strateji yalnızca \CutCS{} için işlediğinden
\Cut{} yerine \CutCS{}'yi seçiyoruz.

\end{document}
